\section{The quadratic mass threshold for a central residual action}
\label{sec:r13-point-residual-mass-threshold}

An ordinary punctured action theorem leaves a current supported
at its central complement.  This note records a quantitative
sufficient condition for removing that residual and the critical
scaling which a singular background can attain.  It concerns the
already defined action of one presentation, as in OL.

\subsection{A fixed-presentation cutoff estimate}

Fix an original ordinary domain, a finite actual coordinate tuple
\(t\), an original compact complementary smooth test \(\eta\),
and an allowed original background \(S\).  Retain the quadratic
return identity which writes the relevant action, up to its fixed
sign, as
\[
                  A=D M_S(Z),\qquad D=d'd''.
\tag{CR1}
\]
The singular full-exact term \(A_0\) in the original decomposition
has not been replaced by a smooth coefficient primitive.  CR1
uses its separately justified scalar-module return identity.

Suppose \(A\) is supported in a relatively closed set \(E\), and
suppose a fixed finite smooth moment map \(t\) has constant value
\(t_0\) on \(E\cap\operatorname{supp}\eta\).  Shrink the fixed
compact test neighborhood so that all cutoffs used below remain
inside the original domain.  Write
\[
 r=|t-t_0|,\qquad
 \rho_\delta=\chi(r/\delta),\qquad
 \chi=1\text{ on }[0,1],\quad \chi=0\text{ on }[2,\infty),
\tag{CR2}
\]
where \(\chi\) is a fixed smooth function, constant near zero.
Although \(r\) need not be smooth at its zero set, \(\rho_\delta\)
is a legitimate original smooth cutoff.  Its value is one on
an ordinary neighborhood of the relevant part of \(E\).
Neither injectivity of \(t\) nor exhaustion of \(E\) by its
zero fiber is needed for the estimate.

Assume the retained degree-one positive-module order estimate
for the represented bounded tensor \(Z\), and fix a positive
control \(Q\) with \(\pm[Z]\le C Q\) on this presentation.
Choose one positive complementary dominator \(\Xi_\eta\) for
the finitely many coefficient forms obtained from \(\eta\),
\(d'\eta\), \(d''\eta\), and \(D\eta\) by adjoining the
coordinate differentials of \(t\).  This is the ordinary finite
coefficient domination used in the retained order interface.
In particular it is chosen after the fixed test and before
\(\delta\).  Let
\[
                    \mu_{S,\eta}=Q\wedge S\wedge\Xi_\eta
\tag{CR3}
\]
denote its finite positive top-degree measure on the compact
test neighborhood.  The existence and positivity of this
measure are part of the displayed module hypotheses, rather
than an identification of the different positivity cones.

\begin{proposition}[Second-order residual bound]
\label{prop:r13-quadratic-residual-mass-threshold}
For \(0<\delta<1\),
\[
 |\langle A,\eta\rangle|
       \le C_\eta\delta^{-2}
            \mu_{S,\eta}\bigl(\operatorname{supp}\eta
                                  \cap\{r\le2\delta\}\bigr).
\tag{CR4}
\]
Consequently the sufficient condition
\[
  \mu_{S,\eta}\bigl(\operatorname{supp}\eta
                      \cap\{r\le2\delta\}\bigr)=o(\delta^2)
\tag{CR5}
\]
removes this residual pairing.  If CR5 holds for every allowed
original \(S\) and every original test, in their actual finite
presentations, then the punctured action identity extends
across \(E\) with those same quantifiers.
\end{proposition}

\begin{proof}
Ordinary locality and the support condition give
\[
 \langle A,\eta\rangle
   =\langle A,\rho_\delta\eta\rangle
   =\langle M_S(Z),D(\rho_\delta\eta)\rangle,
\tag{CR6}
\]
with the harmless fixed current sign suppressed.  The full
Leibniz rule contains the four types
\[
 \rho_\delta D\eta,\qquad
 d'\rho_\delta\wedge d''\eta,\qquad
 d''\rho_\delta\wedge d'\eta,\qquad
 D\rho_\delta\wedge\eta.
\tag{CR7}
\]
Their coefficient bounds in the fixed coordinates are,
respectively, \(C_\eta\), \(C_\eta\delta^{-1}\),
\(C_\eta\delta^{-1}\), and \(C_\eta\delta^{-2}\).
All are supported in \(\operatorname{supp}\eta\cap\{r\le2\delta\}\).
The finite positive domination and \(\pm[Z]\le C Q\)
therefore give CR4.  This proof includes the mixed terms
in CR7.  It does not restrict a current to a radial line.
Taking \(\delta\downarrow0\) under CR5 gives the stated
conclusions.  Each later actual refinement merely gives
another legal instance of the same original order estimate;
no new test class is inferred.
\end{proof}

The assertion is conditional on CR5.  A finite radial fiber
may contain additional points away from \(E\); their mass
also appears in CR4.  The estimate does not discard them
using the vanishing of the action there.  Proving a better
carrier estimate, or using a different family of original
cutoffs, is a separate way to improve CR4.

\subsection{A closed positive background can attain the critical rate}

The rate in CR5 does not follow solely from bounded finite-PS
coefficients, separate closedness and positivity of \(S\),
and point support of \(D M_S(Z)\).  The following supplied
real example retains all of those latter properties.

Use real superforms on \(\mathbf R^3\) with coordinates
\((x,y,z)\).  Let \(H=\{z=0\}\), let \(S=[H]\) be its
positive integration current of bidegree \((1,1)\), and set
\[
 Z=\tfrac12{\bf1}_{\{x>0,y>0\}}
                   (d'x\wedge d''y+d'y\wedge d''x).
\tag{CR8}
\]
Give all coefficient faces their tangential pullbacks.  This
is a bounded finite closed-cell PS tensor with compatible
face pullbacks: on a face contained in \(x=0\) or \(y=0\)
both terms pull back to zero.  The full plane \(H\) is
separately \(d'\)- and \(d''\)-closed, and \([H]\) is
positive for all the usual positivity tests in this degree.
The facewise product in this example is simply
\(M_S(Z)=\iota_*[\iota^*Z]\), where \(\iota:H\hookrightarrow
\mathbf R^3\).

Normalize the positive supervolume on \(H\) by
\(d'x\wedge d''x\wedge d'y\wedge d''y\).  Direct
distributional differentiation gives
\[
               D M_S(Z)=-\delta_0.
\tag{CR9}
\]
Indeed the cross entry of the symmetric coefficient matrix
is \(Z_{xy}=\tfrac12{\bf1}_{x>0}{\bf1}_{y>0}\), while
its two diagonal entries vanish.  In these conventions the
coefficient of \(D[Z]\) on the plane is
\(\partial_y^2 Z_{xx}+\partial_x^2Z_{yy}
             -2\partial_x\partial_y Z_{xy}\), which is
\(-\delta_{(0,0)}\).  The proper linear inclusion commutes
with these distributional derivatives, proving CR9.

For the ordinary Euclidean positive coefficient control,
the measure CR3 on \(H\) is a positive constant multiple
of its two-dimensional area.  Therefore
\[
                   \mu_{S,\eta}(B_\delta(0))\asymp\delta^2
\tag{CR10}
\]
when the fixed test is nonzero near the origin.  The
second-derivative factor in CR4 is exactly critical.

This example does not come from an actual analytic-zero
relation.  Its ambient \(D[Z]\) is nonzero on the
\(z\)-axis, so it does not satisfy the original ambient
closed-current input either.  Its conclusion is exactly
that bounded coefficients and a closed positive background
do not supply a general point-support removal theorem.
The missing original-range input must use the actual-zero
structure, a stronger background mass estimate, or an exact
residual pairing argument.  The local positive-face and
punctured identities alone do not provide that input.

\paragraph{Sources and scope.}
The analytic premises are the retained QA scalar-module
return and the fixed-presentation positive order interface,
as used in NSA12, together with ordinary locality OL3.
CR4 and the real computation CR8--CR10 are the new elementary
deductions.  No support theorem for general currents, no
ordinary-to-GK passage, and no analytic realization of the
real example is asserted.
